\Needspace{8\baselineskip}
\section{Challenges and future directions}
\label{sec:future}

\input{tables/tab_practice}

Nine problems follow from the findings of \S\ref{sec:findings} and the evaluation gaps of
\S\ref{sec:evaluation}; several are the medical instance of open problems the general-domain field has
stated for itself~\cite{sharkey2025open}, and the first is a problem for evidence synthesis in this area
rather than for modelling, which constrains everything after it.

\textit{(0) An atomic unit for representation-level claims.} The record of \S\ref{sec:findings} (study
$\times$ family $\times$ evidence type $\times$ locus) is the finest unit the reviewed papers report
consistently, and it is not atomic, which is why the statements name disagreeing records rather than
count directions. What is needed is a reporting convention: a paper that reports a representation-level result should state, per concept,
model, dataset and end-point, the quantity, the comparison it is read against, and the validity test run on
\emph{that} outcome. A sheet in that form, released beside the paper, would let a later synthesis count
outcomes rather than records. The same convention needs a probe control that medical images do not yet
have. A control task in the sense of Hewitt and Liang~\cite{hewitt2019designing} assigns each input
\emph{type} a fixed random label, so that a probe with capacity but no selectivity scores as well on the
control as on the task; language has recurring types and medical images usually do not, and under a
patient-level split a type has to cross the split --- scanner, site, protocol, view or an annotated patch
class --- or the control measures nothing. Supplement S8c sets out the three designs that should be kept
apart (chance calibration, a control task proper, a memorisation test); none is a citation one can simply
apply, so the selectivity of a medical probe is at present unreported rather than poor.

\textit{(1) Concept preservation across the connector.} Medical multimodal LLMs pass visual evidence through a
connector, where general-domain work finds distinct token geometry and uneven language-model use
(\S\ref{sec:tracing}). One study probes class labels at encoder, connector and language model in fourteen
models, finding fidelity loss~\cite{zhu2026lost}; none compares named-concept decodability or
controllability across it. A solution should measure identical concepts on both sides and test whether
upstream interventions reach reports: probe the same finding across RAD-DINO features, MAIRA-2
projector output and every language-model layer on MIMIC-CXR.

\textit{(2) From decodable to used: paired evidence and validity standards.} High within-study decoding
accuracy is reported throughout (\S\ref{sec:decoding}); interventions are input-dependent, model-specific or null
(\S\ref{sec:intervention}). Only frozen-encoder designs pair probes with residualisation or channel
knockout~\cite{kim2026encoding,nooralahzadeh2026sparse}; no multimodal LLM study combines same-locus probing
and controlled steering with one behavioural end-point, though random directions cause sizeable
changes~\cite{korznikov2025the} and steering vectors are
unidentifiable~\cite{venkatesh2026nonident}. We therefore propose a screening diagnostic for direction-selective sensitivity, an authors'
proposal awaiting benchmarking rather than a validated standard: the steering effect along a direction
$v_c$ built for concept $c$ (\S\ref{sec:methods}, Eq.~(\ref{eq:effect})) should exceed the 95th
percentile of the effects of $K$ pre-specified control directions matched to $v_c$ in locus, norm and
activation scale, each averaged over held-out inputs disjoint from those that built $v_c$, with the
randomisation $p$-value carried only by the label-permuted control family refit through the same
pipeline. The construction, the choice of $K$, the magnitude--sign grid, the patient-level resampling
and a worked example are in supplement S8b; anything short of a pre-registered primary cell or a
grid-maximum statistic is an exploratory screen and should be labelled one.

Supplement S8b lists the fields a steering design should report; the further checks a concept-specific
causal contribution requires are set out in \S\ref{sec:methods}, and no reviewed medical study reports them
all. The paired probe-and-steer design is the natural first test, and the connector experiment of item (1)
its natural setting.

\textit{(3) Reproducible discovery.} SAE features and factorised directions are unstable: absorption and
splitting are structural~\cite{chanin2024a} and only a minority recur across
seeds~\cite{paulo2025sparse,sainsbury2026sparse}. A feature that does not survive retraining cannot be an audit identifier, so progress needs cross-seed
matching, name validation and stable identifiers; universal and concept-anchored dictionaries and crosscoders
across seeds or checkpoints are the starting
points~\cite{thasarathan2025universal,nasirisarvi2025sparc,minder2025overcoming}. An experiment could train two
differently seeded dictionaries on UNI or Virchow features, report the matched-feature fraction and validate
names against CAMELYON regions.

\textit{(4) Shared concept resources and evaluation standards.} Concept banks are rarely clinician-validated
(\S\ref{sec:eval-what}), and protocols vary by dataset, split, label, decoder, width, sparsity and end-point
(\S\ref{sec:eval-compare}). A shared benchmark should define ontology-aligned vocabularies with inter-rater agreement, cover
confounders, shortcuts and findings, fix splits and decoders across public encoders and multimodal LLMs, and
score faithfulness, stability and specificity separately, reporting selectivity rather than accuracy; without
new annotation a first version could probe RAD-DINO and MedGemma with VinDr-CXR and PadChest labels, UNI and
CONCH with CAMELYON regions and MONET with SkinCon. Two reporting cards would make results commensurable, one for
decoding and one for discovery, fixing decoder, split and locus on one side and dictionary, naming,
validation, stability and reconstruction on the other; supplementary Table~S2 and
Table~S16 give the fields.

\textit{(5) Interpretability of adaptation.} Medical foundation models are rarely used as released, and the
evidence on what adaptation changes is thin (\S\ref{sec:geometry}), with checkpoint similarity and cross-model
dictionaries only beginning to reach
medicine~\cite{minder2025overcoming,ahmed2026medseglens}. Which directions survive site transfer, and what
low-rank updates create or destroy, matters for deployment monitoring: an experiment could adapt a public
encoder to a new site, report per-layer similarity and probe drift, and train a crosscoder to name the changes.

\textit{(6) Prospective human and regulatory evaluation.} Medical-device guidance requires declared inputs and
outputs, performance measures, bias sources and traceable
documentation~\cite{fda2025aiguidance,euaiact2024}. Feature-space methods could document encoded attributes
for bias assessment, but no reviewed study links such a measurement to a decision and reader studies are rare
(\S\ref{sec:eval-bench}). A trial could show or withhold a shortcut audit from developers choosing
encoders, recording their choices and the deployed models' subgroup performance.

\textit{(7) Aggregator, volumetric and multi-image loci.} The loci of \S\ref{sec:loci} follow the
two-dimensional single-image path, yet whole-slide aggregators supply the commonest pathology explanations
(\S\ref{sec:tracing}), volumetric encoders pool sub-volumes before probing, and multi-image report generators
encode time at positions unseen by the encoder. None of the three is separately controlled, which is where the
sublocus regularity of \S\ref{sec:findings} bites: future coding should add an aggregator locus and a pooling
declaration, and test aggregator attention against occlusion, attention-MIL weights on UNI or TITAN features
against CAMELYON tiles~\cite{sadanandan2026attention}.

\textit{(8) The text side and the linguistic prior.} The text encoder is one of the two loci with almost no
single-locus core claim (\S\ref{sec:taxonomy}), and the language prior remains unmeasured after fusion. Decoder
dictionaries can reveal report templates, and paraphrase sensitivity is measurable as text-encoder geometry
rather than zero-shot accuracy. An experiment could feed blank or shuffled images to a public CheXpert report
generator and report, per finding, the decodability that survives image removal.

None of the eight experiments needs new instruments: open suites, closed-form linear
erasure~\cite{belrose2023leace}, released dictionaries~\cite{abdulaal2024an,nooralahzadeh2026universal},
the expert-verified scorers of \S\ref{sec:eval-data}~\cite{kim2024transparent} and the models of
Table~\ref{tab:models} suffice. Table~\ref{tab:practice} maps the seven recurring practitioner questions to
a locus, family, validity and transport controls and licensed claim, and supplementary Table~S12 assembles
them into a minimum audit of one attribute, with the inputs, splits, nulls and stopping criteria of each
step. What is feasible depends on access: embeddings alone support permutation- and chance-calibrated probing on a
patient-level split, geometry and site measurement, and erasure before a downstream evaluation; open
encoder weights add a randomly initialised baseline and layer- and pooling-resolved probes, so a sublocus
can be named; and only an open encoder--connector--language-model stack supports tracing, patching and the
direction-selectivity screen, the one tier that licenses a statement about use.
